\documentclass[10pt,a4paper]{amsart}
\usepackage[margin=19mm]{geometry}
\usepackage{amsmath,amssymb,mathtools}
\usepackage[scr=rsfs,cal=euler]{mathalfa}
\usepackage{xfrac}
\usepackage[unicode,hidelinks,bookmarks=false]{hyperref}
\newcommand{\ie}{i.e.\ }
\newcommand{\cf}{cf.\ }
\newcommand{\resp}{respectively\ }
\newcommand{\Subsection}{\subsection}
\providecommand{\nobreakdash}{\nobreak\hbox{-}}
\setlength{\emergencystretch}{2em}
\multlinegap0pt
\usepackage{mathtools}
\usepackage{amssymb}
\usepackage[shortlabels]{enumitem}
\usepackage{xcolor}
\newtheorem{theorem}{Theorem}
\newtheorem{proposition}{Proposition}
\newtheorem{lemma}{Lemma}
\newcommand{\C}{\mathbb C}
\newcommand{\E}{\mathbb E}
\renewcommand{\P}{\mathbb P}
\newcommand{\Pp}{\mathbb P}
\newcommand{\cir}{\mathrm{circ}}
\title{Circular Laws for Cyclic Full-Block Bands:\\ Discrete and Bounded-Density Regimes}
\author{Yi Han}
\date{Source: arXiv:2609.01295v2 (CC BY 4.0)}
\begin{document}
\maketitle

\noindent\emph{Source and licence.} Circular Laws for Cyclic Full-Block Bands:\\ Discrete and Bounded-Density Regimes. Version arXiv:2609.01295v2 (CC BY 4.0), distributed by arXiv under the Creative Commons Attribution 4.0 International licence, http://creativecommons.org/licenses/by/4.0/, as recorded in the arXiv licence metadata for this exact version and shown on the abstract page https://arxiv.org/abs/2609.01295v2. Full licence notice: \texttt{licenses/arxiv-2609.01295-CC-BY-4.0.txt}.\par\smallskip

\noindent\textit{Editorial scope.} Counted unit: the article's proposition prop:local-terminal (source0.tex lines 3179--3214). The article's complete original proof of it is delivered with it (source0.tex lines 5000--5402, one \texttt{proof} environment of 403 lines, which the article places away from the statement). No mathematics is written, repaired or re-derived: the statement and the proof are the article's own lines, in the article's own notation, and no retained line is substituted or re-flowed.  Retained uncounted as the source-local context the proof needs: lines 626--647; lines 649--675; lines 3072--3079; lines 3092--3104; lines 4789--4825; lines 4935--4955; lines 4992--4997.  Nothing was omitted from any retained span, so the package carries no omission record.  No retained line contains a citation, so the wrapper carries no bibliography and no bibliography entry.  No retained line contains a figure environment, an image-inclusion command or drawing code, so no image is delivered and the package declares no asset dependency.  Editorial formatting only, in addition: an AMS \texttt{amsart} preamble that restates the article's own declarations and macros used by the retained lines, namely the environments \texttt{theorem}, \texttt{proposition}, \texttt{lemma} on separate counters, together with the standard packages those lines need.  The retained environments print in this document as Theorem 1, Proposition 1, Lemma 1 in order of appearance, whereas the article prints the same items with the numbering of its own full document.  The complete native source of the article as distributed in the arXiv e-print for this exact version is bundled unchanged as \texttt{sources/209093/source0.tex}.
Define the variance-normalized blocks
\[
 \mathsf A_j=(3W)^{-1/2}A_j,\qquad
 \mathsf B_j=(3W)^{-1/2}B_j,\qquad
 \mathsf C_j=(3W)^{-1/2}C_j
\]
and the cyclic block-tridiagonal matrix
\begin{equation}
 X_N=
 \begin{pmatrix}
  \mathsf A_1&\mathsf B_1&&&\mathsf C_1\\
  \mathsf C_2&\mathsf A_2&\mathsf B_2&&\\
  &\ddots&\ddots&\ddots&\\
  &&\mathsf C_{m-1}&\mathsf A_{m-1}&\mathsf B_{m-1}\\
  \mathsf B_m&&&\mathsf C_m&\mathsf A_m
 \end{pmatrix}.
 \label{eq:model}
\end{equation}
Thus row \(j\) of the block equation contains
\(\mathsf C_j\) at the left neighbor, \(\mathsf A_j\) at the current
site, and \(\mathsf B_j\) at the right neighbor.  Every scalar row has
exactly \(3W\) independent entries of variance \((3W)^{-1}\).

\begin{theorem}[Subgaussian, possibly discrete, atoms]
\label{thm:main}
Let \(X_N\) be the periodic full-block band matrix defined in
\eqref{eq:model}.  Assume that \(\xi\) is real and
\(\|\xi\|_{\psi_2}\le K\), where \(K\) is fixed.  Here
\[
 \|Y\|_{\psi_2}:=\inf\left\{t>0:
   \E\exp\left(\frac{|Y|^2}{t^2}\right)\leq2\right\}
\]
denotes the subgaussian Orlicz norm.  Assume further that
\begin{equation}
 W=W_N\longrightarrow\infty,
 \qquad \frac{W}{\log N}\longrightarrow\infty.
 \label{eq:bandwidth}
\end{equation}
Then the empirical eigenvalue measure of \(X_N\) converges in probability
to the uniform probability measure \(\mu_{\cir}\) on the unit disk.
More precisely, for every fixed \(z\in\C\),
\begin{equation}
 \frac1N\log|\det(X_N-zI_N)|
 \xrightarrow{\mathbb P}
 U_{\cir}(z)
 \label{eq:logdet-goal}
\end{equation}
In particular, \(\Pp\{\det X_N=0\}\to0\), so \(X_N\) is nonsingular
with high probability.
\end{theorem}

\begin{align}
 C_L\psi_-+(A_L-zI_W)\psi_L+B_L\psi_C&=0,
 \label{eq:local-packet-left}\\
 C_C\psi_L+(A_C-zI_W)\psi_C+B_C\psi_R&=0,
 \label{eq:local-packet-centre}\\
 C_R\psi_C+(A_R-zI_W)\psi_R+B_R\psi_+&=0.
 \label{eq:local-packet-right}
\end{align}

With both rows and columns ordered as \((L,C,R)\), the packet mask related to \eqref{eq:local-packet-left}-\eqref{eq:local-packet-right} is
\begin{equation}
 M_{\partial}=
 \begin{pmatrix}
  1&1&0\\
  1&1&1\\
  0&1&1
 \end{pmatrix},
 \label{eq:local-mask}
\end{equation}
where each \(1\) denotes a complete \(W\times W\) random block.  Equivalently,
the only missing rectangles are \(L_{\rm row}\times R_{\rm col}\) and
\(R_{\rm row}\times L_{\rm col}\).

\begin{lemma}[Two-sided coordinate skeleton]
\label{lem:local-rrqr}
Let \(Q\in\C^{n\times n}\), let \(\tau\geq1\), and write
\(r:=\#\{j:s_j(Q)>\tau\}\).  There are row and column sets
\(I,J\subset[n]\), each of cardinality \(r\), such that
\begin{equation}
 K_{\mathrm{piv}}:=Q_{I,J}
 \label{eq:local-Kpiv}
\end{equation}
is invertible when \(r>0\).  Define
\begin{equation}
 X_{\mathrm{skel}}:=K_{\mathrm{piv}}^{-1}Q_{I,J^c},\quad
 Y_{\mathrm{skel}}:=Q_{I^c,J}K_{\mathrm{piv}}^{-1},\quad
 E_0:=Q_{I^c,J^c}-Y_{\mathrm{skel}}K_{\mathrm{piv}}X_{\mathrm{skel}}.
 \label{eq:local-skeleton-definitions}
\end{equation}
After a row-column permutation applied to $Q$, we can write
\begin{equation}
 Q=
 \begin{pmatrix}
  K_{\mathrm{piv}}&K_{\mathrm{piv}}X_{\mathrm{skel}}\\
  Y_{\mathrm{skel}}K_{\mathrm{piv}}&
  Y_{\mathrm{skel}}K_{\mathrm{piv}}X_{\mathrm{skel}}+E_0
 \end{pmatrix},
 \label{eq:local-skeleton-identity}
\end{equation}
and, with the universal choice \(C_{\rm RRQR}=4\), for all large $n$,
\begin{equation}
 s_j(K_{\mathrm{piv}})\geq n^{-C_{\rm RRQR}}s_j(Q),\quad
 \|X_{\mathrm{skel}}\|+\|Y_{\mathrm{skel}}\|
 \leq n^{C_{\rm RRQR}},\quad
 \|E_0\|\leq n^{C_{\rm RRQR}}\tau.
 \label{eq:local-skeleton-bounds}
\end{equation}
The first bound holds for \(j\leq r\).  For \(r=0\), the pivot is empty
and \eqref{eq:local-skeleton-identity} means \(Q=E_0\).
\end{lemma}

\begin{lemma}[The deformed-square input used below]
\label{lem:local-cook-input}
Fix constants \(L<\infty\) and \(0<c_*\leq C_*<\infty\).  Let
\[
 Z_n=(a_{ij}\xi_{ij})_{i,j\leq n},
 \qquad c_*\leq a_{ij}\leq C_*,
\]
where the \(\xi_{ij}\) are independent copies of the atom in
Theorem~\ref{thm:main}.  If \(D_n\in\C^{n\times n}\) is deterministic
and \(\|D_n\|\leq n^L\), then there are constants
\(\beta_L,C_L,c_L>0\), depending only on
$L,K,c_*,C_*$, such that
\begin{equation}
 \P\{s_{\min}(Z_n+D_n)\leq n^{-\beta_L}\}
 \leq C_L\sqrt{\frac{\log n}{n}}+e^{-c_Ln}.
 \label{eq:local-cook-input}
\end{equation}
The same bound holds conditionally when \(D_n\) is measurable with
respect to an independent sigma-field and satisfies
$\|D_n\|\leq n^L$ almost surely.
\end{lemma}

\begin{equation}
 p_{Q_O}=\sigma_W^{-3W}\widetilde p_{\widetilde Q_O},\qquad
 \mathfrak C(Q_O)
 =\sigma_W^{-3W}\widetilde{\mathfrak C}(\widetilde Q_O).
 \label{eq:local-terminal-scaling}
\end{equation}

\begin{proposition}[Three-block terminal estimate]
\label{prop:local-terminal}
Assume \(|\zeta_W|=\sqrt{3W}|z|\leq W^{K_z}\) for a fixed \(K_z\).  Then uniformly in
\(Q_O\in\C^{2W\times2W}\), for every \(T>0\),
\begin{equation}
 \E\min\left\{T,
  \log_+\frac{\mathfrak C(Q_O)}
                   {|p_{Q_O}(\Xi)|}\right\}
 \leq CW\log W+p_WT,
 \qquad
 p_W\leq C\sqrt{\frac{\log W}{W}}+e^{-cW}.
 \label{eq:local-terminal-capped}
\end{equation}
The capped loss is \(T\) at a zero of the determinant.  In particular,
\begin{equation}
 \P\{p_{Q_O}(\Xi)=0\}\leq p_W.
 \label{eq:local-terminal-zero}
\end{equation}
On the event
\begin{equation}
 \mathcal E_{\max}:=\left\{\max_{e\in\mathcal I_{\partial}}|\xi_e|
                 \leq C_0\sqrt W\right\},
 \qquad \P(\mathcal E_{\max}^c)\leq e^{-cW},
 \label{eq:local-max-event}
\end{equation}
where \(C_0=C_0(K)\) is sufficiently large,
one also has
\begin{equation}
 \log_+\frac{|p_{Q_O}(\Xi)|}
                  {\mathfrak C(Q_O)}
 \leq CW\log W.
 \label{eq:local-terminal-reverse}
\end{equation}
All assertions remain valid conditionally when \(Q_O\) is measurable with
respect to a sigma-field independent of the seven packet blocks.
\end{proposition}

\begin{proof}[Proof of Proposition~\ref{prop:local-terminal}]
By \eqref{eq:local-terminal-scaling}, the two logarithmic ratios and the
zero event are invariant under the scaling by $\sigma_W$.  Since
\(Q_O\mapsto\sigma_WQ_O\) is a bijection of
\(\C^{2W\times2W}\), it therefore suffices to prove the asserted bounds
for \(\widetilde p_{Q_O}\) and
\(\widetilde{\mathfrak C}(Q_O)\), with \(Q_O\) now denoting the
row-scaled deformation for the duration of this proof.


\textbf{Step 1: removal of $Q_O$ via linear algebra.}
Choose a large fixed \(K_0>K_z+C_{\rm RRQR}+10\), define
\(\tau_W:=W^{K_0}\), and
set
\begin{equation}
 r:=\#\{j:s_j(Q_O)>\tau_W\},\qquad n_{\rm res}:=3W-r.
 \label{eq:local-r-nres}
\end{equation}
Thus \(r\leq2W\) and \(n_{\rm res}\geq W\).  Apply
Lemma~\ref{lem:local-rrqr} to \(Q_O\), and extend its coordinate
decomposition by a zero block on the central coordinates.  We obtain
sets \(I,J\subset O\), \(|I|=|J|=r\), and an exact decomposition of the
deterministic matrix \(\operatorname{Emb}_O(Q_O)\) of the form
\eqref{eq:local-skeleton-identity}.  The rows of \(Y_{\mathrm{skel}}\) and
the columns of \(X_{\mathrm{skel}}\) indexed by \(C\) are zero.
When \(r=0\), the residual below is simply the
full matrix with a polynomial-norm deformation and the pivot matrix does not exist.  We write the remaining
formulas for \(r>0\), with this empty-pivot convention understood.

We absorb the deterministic shift into the row-scaled packet matrix:
\begin{equation}
 \Delta_z:=\widetilde\Delta+\zeta_WI_{3W}.
 \label{eq:local-Delta-tilde}
\end{equation}
Then we split \(\Delta_z\) in the same coordinates as $\Delta$:
\begin{equation}
 \Delta_z=
 \begin{pmatrix}(\Delta_z)_{11}&(\Delta_z)_{12}\\
                 (\Delta_z)_{21}&(\Delta_z)_{22}\end{pmatrix},
 \qquad
 K_\Delta:=K_{\mathrm{piv}}+(\Delta_z)_{11}.
 \label{eq:local-Delta-split}
\end{equation}
This is the pivot--residual split: \((\Delta_z)_{11}\) is the
\(r\times r\) pivot perturbation, whereas \((\Delta_z)_{22}\) is the
entire \(n_{\rm res}\times n_{\rm res}\) residual.  They are not the two
complete iid squares used below; those squares, denoted by \(Z_1,Z_2\),
will be selected inside the residual in \eqref{eq:local-two-cook-blocks}.

Let \(\mathcal F_{\rm piv}\) be generated by the packet entries occurring
in \((\Delta_z)_{11},(\Delta_z)_{12},(\Delta_z)_{21}\), and define
\begin{equation}
 \mathcal E_{\rm piv}:=
 \left\{\max_{(i,j)\in\{(1,1),(1,2),(2,1)\}}
 \|(\Delta_z)_{ij}\|\leq W^{C_\Delta}\right\},
 \qquad C_\Delta:=\max\{K_z,3\}+1.
 \label{eq:local-pivot-exposure-event}
\end{equation}
The standard subgaussian norm bound gives
\(\P(\mathcal E_{\rm piv}^c)\leq e^{-cW}\).  These variables lie outside
the residual, so \(\mathcal E_{\rm piv}\) is independent of the two
complete residual squares chosen below.  We use it only as a good event;
we do not condition those squares on it.  On \(\mathcal E_{\rm piv}\),
the choice of \(K_0\) gives
\(\|K_{\mathrm{piv}}^{-1}(\Delta_z)_{11}\|\leq1/2\), and hence
\begin{equation}
 |\det K_\Delta|
 \geq e^{-CW\log W}\prod_{j=1}^r s_j(Q_O).
 \label{eq:local-pivot-det}
\end{equation}
Indeed,
 \(\det K_\Delta=\det K_{\mathrm{piv}}
 \det(I+K_{\mathrm{piv}}^{-1}(\Delta_z)_{11})\);
the second determinant has modulus at least \(2^{-r}\), while strong
RRQR bounds the first by
\(W^{-C_{\rm RRQR}r}\prod_{j\leq r}s_j(Q_O)\).

To change the matrix to a further diagonal form, we denote by
\begin{align}
 G_{21}&:=(\Delta_z)_{21}
          -Y_{\mathrm{skel}}(\Delta_z)_{11},
 &G_{12}&:=(\Delta_z)_{12}
          -(\Delta_z)_{11}X_{\mathrm{skel}},
 \label{eq:local-Gs}\\
 F&:=E_0-Y_{\mathrm{skel}}(\Delta_z)_{12}
       -(\Delta_z)_{21}X_{\mathrm{skel}}\notag\\
  &\quad+Y_{\mathrm{skel}}(\Delta_z)_{11}X_{\mathrm{skel}}
       -G_{21}K_\Delta^{-1}G_{12}.
 \label{eq:local-F}
\end{align}
Here and below, an inverse is extended by zero on its singular set; this
is a Borel matrix-valued map.  Identities involving that inverse are used
only on the good event on which it is a genuine inverse.  In particular,
on \(\mathcal E_{\rm piv}\), direct multiplication gives the exact
factorization
\begin{align}
 \operatorname{Emb}_O(Q_O)+\Delta_z
 &={}
 \begin{pmatrix}
  I&0\\
  (Y_{\mathrm{skel}}K_{\mathrm{piv}}
    +(\Delta_z)_{21})K_\Delta^{-1}&I
 \end{pmatrix}
 \begin{pmatrix}K_\Delta&0\\0&(\Delta_z)_{22}+F\end{pmatrix}
 \notag\\
 &\quad\times
 \begin{pmatrix}
  I&K_\Delta^{-1}(K_{\mathrm{piv}}X_{\mathrm{skel}}
                    +(\Delta_z)_{12})\\
  0&I
 \end{pmatrix}.
 \label{eq:local-CUR}
\end{align}
The first
and third factors are determinant-one shears.  More importantly, every
unbounded occurrence of \(K_{\mathrm{piv}}\) cancels from \(F\).
Here is the promised norm check.  On \(\mathcal E_{\rm piv}\), the three
blocks of \(\Delta_z\) occurring in \eqref{eq:local-F} have operator norm
at most \(W^{C_\Delta}\).  The RRQR bounds (Lemma \ref{lem:local-rrqr}) and the
choice of \(K_0\) give, for fixed exponents independent of \(W\),
\begin{align}
 \|X_{\mathrm{skel}}\|+\|Y_{\mathrm{skel}}\|
 &\leq W^{C_{\rm RRQR}},&
 \|E_0\|&\leq W^{K_0+C_{\rm RRQR}},\notag\\
 \|K_\Delta^{-1}\|
 &\leq2\|K_{\mathrm{piv}}^{-1}\|
 \leq W^{C_{\rm RRQR}-K_0},&
 \|G_{12}\|+\|G_{21}\|
 &\leq W^{C_{\rm RRQR}+C_\Delta+1}.
 \label{eq:local-CUR-term-bounds}
\end{align}
Substituting these four estimates into the five terms of
\eqref{eq:local-F} shows, term by term, that on
\(\mathcal E_{\rm piv}\)
\begin{equation}
 \|F\|\leq W^{K_1}
 \label{eq:local-F-bound}
\end{equation}
for a fixed \(K_1=K_1(K_0,C_{\rm RRQR})\).  This eventwise bound will
show below that the global norm truncations are
inactive on \(\mathcal E_{\rm exp}\).

\textbf{Step 2: analyzing the remaining block via polynomial LSV.}
We first describe the structure of the matrix $(\Delta_z)_{22}+F$ in the original coordinates.
Let \(L_{\rm r},R_{\rm r}\) be the remaining outer row sets and
\(L_{\rm c},R_{\rm c}\) the remaining outer column sets, after removal of labels $I$ and $J$.  Denote by
\begin{equation}
 a=|L_{\rm r}|,\quad b=|R_{\rm r}|,\quad
 c=|L_{\rm c}|,\quad e=|R_{\rm c}|.
\end{equation}
Since equally many outer rows and columns were removed (no central rows
or columns are removed),
\begin{equation}
 a+b=c+e=:s,\qquad n_{\rm res}=W+s.
 \label{eq:local-counts}
\end{equation}
Choose an integer
\begin{equation}
 n_1\in[\max(a,c),W+\min(a,c)]
       \cap[n_{\rm res}/3-1,2n_{\rm res}/3+1].
 \label{eq:local-n1}
\end{equation}
The intersection is nonempty: its left endpoint is at most
\(s\leq2(W+s)/3=2n_{\rm res}/3\), while its right endpoint is at least
\(W\geq(W+s)/3=n_{\rm res}/3\), and integer rounding costs at most one.
Set \(n_2:=n_{\rm res}-n_1\), \(x:=n_1-a\), and \(y:=n_1-c\).  Select \(x\) central rows and \(y\) central
columns.  Together with the remaining left outer coordinates they form a
complete \(n_1\times n_1\) square.  Their complements, together with the
remaining right outer coordinates, form a complete
\(n_2\times n_2\) square.  Neither square meets a forbidden
rectangle of \eqref{eq:local-mask}; the first uses only \(L,C\) coordinates
and the second only \(R,C\) coordinates.  They are row-disjoint,
column-disjoint, and together cover every residual row and column.

Let \(\mathcal F_0\) be generated by all packet entries outside these two
squares.  Conditional on \(\mathcal F_0\), the two retained matrices
\(Z_1,Z_2\) are independent complete iid squares of respective dimensions
\(n_1,n_2\).  On \(\mathcal E_{\rm piv}\), write
\begin{equation}
 (\Delta_z)_{22}+F=
 \begin{pmatrix}
  Z_1+D_{11}&D_{12}\\D_{21}&Z_2+D_{22}
 \end{pmatrix},
 \label{eq:local-two-cook-blocks}
\end{equation}
where every \(D_{ij}\) is \(\mathcal F_0\)-measurable.  Extend the
\(D_{ij}\) measurably off \(\mathcal E_{\rm piv}\), for instance by zero.
By \eqref{eq:local-F-bound} and the subgaussian norm bound for the exposed
cross blocks, there is a fixed \(L_0\) such that the
\(\mathcal F_0\)-measurable event
\begin{equation}
 \mathcal E_{\rm exp}:=\mathcal E_{\rm piv}\cap
 \left\{\max_{i,j\leq2}\|D_{ij}\|\leq W^{L_0}\right\}
 \label{eq:local-full-exposure-event}
\end{equation}
satisfies \(\P(\mathcal E_{\rm exp}^c)\leq e^{-cW}\).

We now make the two conditional applications without conditioning on a
norm event involving \(Z_1\) or \(Z_2\).  For a matrix \(A\), let
\begin{equation}
 \mathcal T_R(A):=A\min\{1,R/\|A\|\},
 \qquad \mathcal T_R(0):=0,
 \label{eq:local-measurable-norm-truncation}
\end{equation}
and set \(\widehat D_{ij}:=\mathcal T_{W^{L_0}}(D_{ij})\).  These
truncated deformations are globally \(\mathcal F_0\)-measurable, have
polynomial norm, and agree with \(D_{ij}\) on
\(\mathcal E_{\rm exp}\).  Lemma~\ref{lem:local-cook-input}, applied
conditionally first to \(Z_1+\widehat D_{11}\), gives, for a fixed
\(\beta_1\), after defining
\[
 B_1:=Z_1+\widehat D_{11},
 \qquad \mathcal G_1:=\{s_{\min}(B_1)\geq W^{-\beta_1}\},
\]
\begin{equation}
 \P(\mathcal G_1^c\mid\mathcal F_0)
 \leq C\sqrt{\frac{\log W}{W}}+e^{-cW}.
 \label{eq:local-cook-one}
\end{equation}
Here both \(n_i\geq n_{\rm res}/3-1\geq W/3-1\), so, for all large
\(W\), \(W\leq4n_i\), while \(n_i\leq3W\).  Thus the estimates furnished
by Lemma~\ref{lem:local-cook-input} at size \(n_i\) imply the displayed
\(W\)-scale estimate: in particular, \(W^{L_0}\leq n_i^{L_0+1}\) for
large \(W\), and
\(e^{-cn_i}\leq e^{-cW/4}\), after which constants are renamed.

To make the second conditioning global, define the Borel truncated inverse
\begin{equation}
 \mathcal R_W(B):=
 \begin{cases}
  B^{-1},&s_{\min}(B)\geq W^{-\beta_1},\\
  0,&s_{\min}(B)<W^{-\beta_1}.
 \end{cases}
 \label{eq:local-truncated-inverse}
\end{equation}
It has norm at most \(W^{\beta_1}\).  After revealing \(Z_1\), define
the globally measurable Schur deformation and its corresponding square by
\begin{equation}
 \widehat S_2:=\widehat D_{22}
 -\widehat D_{21}\mathcal R_W(B_1)\widehat D_{12},
 \qquad \widehat B_2:=Z_2+\widehat S_2.
 \label{eq:local-second-schur}
\end{equation}
Conditional on \(\mathcal F_0\vee\sigma(Z_1)\), its deformation is
globally measurable, has fixed-polynomial norm, and is independent of the
untouched square \(Z_2\).  A second application of
Lemma~\ref{lem:local-cook-input} gives a fixed \(\beta_2\) such that, for
\(\mathcal G_2:=\{s_{\min}(\widehat B_2)\geq W^{-\beta_2}\}\),
\begin{equation}
 \P(\mathcal G_2^c\mid\mathcal F_0\vee\sigma(Z_1))
 \leq C\sqrt{\frac{\log W}{W}}+e^{-cW}.
 \label{eq:local-cook-two}
\end{equation}
To record the determinant identity without defining a Schur complement
only on \(\mathcal G_1\), introduce the artificial bottom block
\begin{equation}
 B_{22}^{\rm art}:=\widehat B_2+\widehat D_{21}
        \mathcal R_W(B_1)\widehat D_{12},
 \qquad
 M_{\rm art}:=
 \begin{pmatrix}B_1&\widehat D_{12}\\
                 \widehat D_{21}&B_{22}^{\rm art}\end{pmatrix}.
 \label{eq:local-artificial-bottom-block}
\end{equation}
On \(\mathcal G_1\), \(\mathcal R_W(B_1)=B_1^{-1}\), and Schur
complementation gives
\[
 \det M_{\rm art}=\det B_1\det\widehat B_2.
\]
On \(\mathcal E_{\rm exp}\), the four norm truncations are inactive;
cancellation in \eqref{eq:local-artificial-bottom-block} then shows that
\(M_{\rm art}\) equals the residual matrix in
\eqref{eq:local-two-cook-blocks}.  Therefore, on
\(\mathcal E_{\rm exp}\cap\mathcal G_1\cap\mathcal G_2\),
\[
 |\det((\Delta_z)_{22}+F)|
 \geq W^{-\beta_1n_1-\beta_2n_2}
 \geq e^{-CW\log W}.
\]
The union bound, \(\P(\mathcal E_{\rm exp}^c)\leq e^{-cW}\), and the two
Cook estimates consequently give, outside probability \(p_W\),
\begin{equation}
 |\det((\Delta_z)_{22}+F)|\geq e^{-CW\log W}.
 \label{eq:local-residual-det}
\end{equation}
Combining \eqref{eq:local-pivot-det}, \eqref{eq:local-CUR}, and
\eqref{eq:local-residual-det} gives
\begin{equation}
 |\widetilde p_{Q_O}(\Xi)|
 \geq e^{-CW\log W}\prod_{j=1}^r s_j(Q_O)
 \label{eq:local-value-lower}
\end{equation}
outside probability \(p_W\).

\textbf{Step 3: from product singular values to the coefficient norm.}
It remains to compare the right side with the coefficient norm.  First
set \(\zeta_W=0\), and denote the corresponding coefficient norm by
\(\widetilde{\mathfrak C}_0(Q_O)\).  An atom monomial is indexed by a
partial permutation \(\pi:I_0\to J_0\).  Its coefficient is
\begin{equation}
 \pm\left(\prod_{i\in I_0}a_{i,\pi(i)}\right)
 \det\bigl(\operatorname{Emb}_O(Q_O)\bigr)_{I_0^c,J_0^c}.
 \label{eq:local-partial-permutation}
\end{equation}
Every nonzero determinant here is a minor of \(Q_O\).  A
\(k\times k\) minor is at most \(\prod_{j\leq k}s_j(Q_O)\).
If \(k<r\), the omitted selected singular values are all larger than
\(\tau_W>1\).  If \(k>r\), the extra singular values are at most
\(\tau_W\) by \eqref{eq:local-r-nres}.  Thus this upper bound extracts only
the very large singular values of $Q_O$; those below
$\tau_W=W^{K_0}$ are absorbed into the polynomial factors.
Thus every coefficient in \eqref{eq:local-partial-permutation} is at
most
\[
 e^{CW\log W}\prod_{j=1}^r s_j(Q_O).
\]
The number of partial permutations of \(3W\) rows and columns is at most
\begin{equation}
 \sum_{t=0}^{3W}\binom{3W}{t}^2t!
 \leq 4^{3W}(3W)!
 \leq e^{CW\log W}.
 \label{eq:local-partial-permutation-count}
\end{equation}
Thus
\begin{equation}
 \widetilde{\mathfrak C}_0(Q_O)
 \leq e^{CW\log W}\prod_{j=1}^r s_j(Q_O).
 \label{eq:local-coeff-upper-zero}
\end{equation}
Conversely, leave precisely the pivot rows (I) and pivot columns (J)
unmatched, and choose a perfect matching in each of the two complete
residual squares constructed in Step~2.  The union of these matchings is
an atom monomial whose coefficient is the product of its deterministic
weights times
\(\det(Q_O)_{I,J}=\det K_{\mathrm{piv}}\).
Strong RRQR in Lemma \ref{lem:local-rrqr} therefore gives
\begin{equation}
 \widetilde{\mathfrak C}_0(Q_O)
 \geq e^{-CW\log W}\prod_{j=1}^r s_j(Q_O).
 \label{eq:local-coeff-lower-zero}
\end{equation}

Now we restore the shift into $\widetilde{\mathfrak C}_0(Q_O)$.  Translating each of the \(3W\) fresh diagonal
variables by \(\zeta_W/a_e\) sends the zero-shift polynomial to the
shifted polynomial.  On the two-dimensional coefficient space of one
multiaffine variable this is a triangular map; its inverse is the
opposite translation.  The norms of the product map and its inverse, generated by this scalar shift, are
at most
\[
 \prod_{e:\,e\text{ is diagonal}}(1+|\zeta_W|/a_e)
 \leq e^{CW\log W}.
\]
Consequently \(\widetilde{\mathfrak C}(Q_O)
=e^{\pm CW\log W}\widetilde{\mathfrak C}_0(Q_O)\), and
\begin{align}
 \widetilde{\mathfrak C}(Q_O)
 &\leq e^{CW\log W}\prod_{j=1}^r s_j(Q_O),
 \label{eq:local-coeff-upper}\\
 \widetilde{\mathfrak C}(Q_O)
 &\geq e^{-CW\log W}\prod_{j=1}^r s_j(Q_O).
 \label{eq:local-coeff-lower}
\end{align}
Equations \eqref{eq:local-value-lower} and
\eqref{eq:local-coeff-upper}, with the bad event charged at height \(T\),
prove \eqref{eq:local-terminal-capped}.  Dividing by \(T\) and letting
\(T\to\infty\) proves \eqref{eq:local-terminal-zero}.

Finally, there are at most \(e^{CW\log W}\) partial-permutation monomials.
On \(\mathcal E_{\max}\), every monomial has modulus at most
\((C_0\sqrt W)^{3W}\).
Cauchy--Schwarz over the coefficient vector gives
\eqref{eq:local-terminal-reverse}.  Centeredness and unit variance also
give the exact Parseval identity
\begin{equation}
 \bigl(\widetilde{\mathfrak C}(Q_O)\bigr)^2
 =\E|\widetilde p_{Q_O}(\Xi)|^2,
 \label{eq:local-parseval}
\end{equation}
since distinct multiaffine monomials are orthonormal in \(L^2\) norm.

By \eqref{eq:local-terminal-scaling}, we may now return to the original
notation $p_q$ and $\mathfrak C(q)$.  The preceding argument treated
deterministic \(Q_O\).  Let \(\mathcal G\)
be a sigma-field such that \(Q_O\) is \(\mathcal G\)-measurable and
\(\Xi\) is independent of \(\mathcal G\), and let \(\mu\) be the law of
\(\Xi\).  The maps \((q,x)\mapsto p_q(x)\) and
\(q\mapsto\mathfrak C(q)\) are Borel.  Thus, with
\[
 \Phi_T(q,x):=\min\left\{T,
   \log_+\frac{\mathfrak C(q)}{|p_q(x)|}\right\},
\]
interpreted as \(T\) when \(p_q(x)=0\), independence and Fubini give
\[
 \E[\Phi_T(Q_O,\Xi)\mid\mathcal G]
 =\int\Phi_T(Q_O,x)\,d\mu(x)
 \leq CW\log W+p_WT
 \quad\text{a.s.}
\]
The same argument with \(\mathbf1_{\{p_q(x)=0\}}\) proves the conditional
zero bound.  The reverse estimate is pointwise in \((q,x)\) on
\(\mathcal E_{\max}\), which is independent of \(\mathcal G\); this
proves the remaining conditional assertions.
\end{proof}

\end{document}
